\documentclass[10pt]{article}
\usepackage[margin=0.6in]{geometry}
\usepackage{amsmath,amssymb,mathtools}
\usepackage{mathptmx}
\DeclareMathAlphabet{\mathcal}{OMS}{cmsy}{m}{n}% keep Computer Modern calligraphic letters
\usepackage{booktabs}
\usepackage{array}
\usepackage{enumitem}
\usepackage{xcolor}
\usepackage{tikz}
\usetikzlibrary{arrows.meta,positioning}
\usepackage[hidelinks]{hyperref}
\usepackage[font=footnotesize,skip=2pt]{caption}
\usepackage{graphicx}
\usepackage{float}
\usepackage[round]{natbib}
\usepackage{titlesec}
\usepackage[normalem]{ulem}
% Workshop 2 change marking: new or changed text in teal, removed text struck out in red
\definecolor{wtwo}{RGB}{0,115,150}
\definecolor{wtwodel}{RGB}{190,30,30}
\newcommand{\add}[1]{\ifmmode{\color{wtwo}#1}\else\leavevmode{\color{wtwo}#1}\fi}
\newcommand{\del}[1]{\ifmmode\text{\color{wtwodel}\sout{$#1$}}\else\leavevmode{\color{wtwodel}\sout{#1}}\fi}
\newenvironment{addblock}{\color{wtwo}}{}
\newcommand{\wone}[1]{{\color{gray}\scriptsize(#1)}}% Workshop 1 value, for comparison
\setlist{nosep,leftmargin=*}
\titlespacing*{\section}{0pt}{5pt}{2pt}
\titlespacing*{\subsection}{0pt}{3pt}{1pt}
\titlespacing*{\paragraph}{0pt}{3pt}{6pt}
\setlength{\parskip}{1pt}
\setlength{\parindent}{0pt}
\setlength{\textfloatsep}{6pt}
\setlength{\floatsep}{6pt}
\setlength{\intextsep}{6pt}

\newcommand{\I}{\mathcal{I}}
\newcommand{\V}{\mathcal{V}}
\newcommand{\E}{\mathcal{E}}
\newcommand{\C}{\mathcal{C}}
\newcommand{\K}{\mathcal{K}}
\newcommand{\Z}{\mathbb{Z}}
\newcommand{\1}{\mathbf{1}}
\newcommand{\B}{\mathsf{B}}
\newcommand{\Rd}{\mathsf{R}}
\newcommand{\As}{\mathrm{A}}
\newcommand{\Pt}{\mathrm{P}}

\title{\vspace{-3.2em}\textbf{Gray-Zone Maritime Control in the First Island Chain:\\ A Networked Escalation Game}\\[2pt]
\large Updates to the Workshop 1 Writeup (working copy for Workshop 2)}
\author{Sean Bowman, Tate Murphy, Matthew Newell\\
\small 6.S974/2.S988/1.S980 Games, Learning, and Security, Fall 2026}
\date{\vspace{-2.9em}}

\begin{document}
\maketitle
\begin{center}\fbox{\parbox{0.95\textwidth}{\small This is a copy of the Workshop 1 report (as submitted, branch
\texttt{workshop1-final}) with the Workshop 2 model changes written in. \add{Text in this color is new or changed.}
\del{Struck-out red text} is removed. Sections not yet touched are unchanged from Workshop 1, including all results,
which will be replaced after the rerun.}}\end{center}


%======================================================================
\section{Setting and research questions}
%======================================================================
\textbf{Contested system:} China and the United States compete, \emph{below the threshold of war}, for day-to-day
control of the waters inside the First Island Chain. China claims most of the South China Sea and contests the
Senkakus, pressing these claims through the persistent presence of coast guard and maritime militia (CCG) hulls
backed by navy (PLAN) ships. The U.S.\ rejects maritime claims beyond what international law allows and keeps the
seas open with a smaller number of destroyers (DDG) on presence and freedom-of-navigation patrols
\citep{helmus2024,green2017}. Every two days, each side decides \emph{where} to station its ships on a network of
zones and \emph{how} to act there: patrol (P) or assert (A: shadowing, close maneuvering, querying and warning, and,
for the U.S.\ only, freedom-of-navigation operations).

\textbf{Why strategic interaction is central:}
control of a zone depends on the ratio of local forces, so each side's best allocation depends on the other's.
Escalation risk is shared. Incidents need both sides present, and war costs both. Neither side sees the other's
entire fleet disposition, only forces present where they also maintain forces.

\textbf{Research questions:}
\begin{enumerate}[label=\textbf{RQ\arabic*.}]
\item \emph{Control vs.\ war risk.} Each side caps how much incident risk it will accept each period (its risk
budget $\varepsilon_i$). How much more control does a side gain by accepting more risk, how much more likely does
war become, and what risk level does each side pick when it knows the other side is picking too?
\item \emph{Game theory vs.\ rules of thumb.} Does a side that plans with game theory (best responses,
equilibrium or maximin plans) do better than one that follows simple rules or moves at random? Can both sides
keep patrolling over many periods because asserting once would end the truce and push both into riskier play from
then on?
\item \emph{Weather and redesign.} Does the plan or Blue redesign that works best in normal weather still work
best across calm, typhoon and extreme weather scenarios, both on average and in the worst scenario?
\end{enumerate}
\textbf{Outcomes (averaged over 200 simulated games per setup):} Blue's value-weighted control share
$\bar c_\B=\frac1\tau\sum_t\sum_{v}V^\B_v c^t_v/\sum_v V^\B_v\in[-1,1]$, war rate and payoffs $u_\B,u_\Rd$. For RQ3,
also expected and worst-case payoff over $\omega$.

\textbf{Changes from the proposal and how we used the feedback:} the model is now a formal game on a graph (\S3).
The best response is an exact MIP whose feasible region is a risk budget, with the proposal's commander types as
budget levels (\S\ref{sec:mip}). The RQs now ask for the budget-game equilibrium, equilibrium play, and robustness
and redesign. We define three Blue redesigns (\S\ref{sec:designs}), propose values for missing data
(Tables~\ref{tab:params}--\ref{tab:zones}), report analytical results and 77 configurations $\times$ 200 seeds
(\S4), and add weather scenarios scored by expected and worst-case payoff, as Prof.\ Sapsis asked (\S\ref{sec:state},
\S4.4).

%======================================================================
\section{Literature and current practice}
%======================================================================
\textbf{Documented practice:} 
\citet{helmus2024} describe China's use of coast guard and maritime militia vessels to 
maintain presence and obstruct maritime activities while seeking to avoid a conventional military response. 
\citet{green2017} emphasize credible commitments, coordinated allied responses, and acceptance of calculated escalation risk while exercising restraint.
U.S.\ practice is routine presence and freedom-of-navigation patrols, with encounters managed
by the Code for Unplanned Encounters at Sea \citep{cues2014}. From these sources we take the actions,
the asymmetric hull types, and the idea that incidents, not shots, drive escalation.

\textbf{Models:} 
\citet{wasser2019} developed a structured gray-zone game with explicit objectives and probabilistic adjudication. We automate the rules to support repeated simulation and strategy comparison. \citet{shieh2012} document a deployed Stackelberg model of randomized Coast Guard port patrols. 
\citet{washburn1995} model network interdiction as a two-person zero-sum game. \citet{schelling1966} represents escalation as a shared,
manipulated risk, motivating our incident hazard.

\textbf{What we add:} we combine maritime allocation, evolving control, and shared escalation risk in a general-sum game where both sides contest control simultaneously and war costs both. Computational rules support comparisons of strategies and designs across thousands of runs. Numerical parameters are our
assumptions (Table~\ref{tab:params}), not values derived from documented literature.

%======================================================================
\section{Game formulation}
%======================================================================
\textbf{One period $t$ (timeline):} (1) Nature draws the sea state $w^t_r$ of each region (scenario $\omega$).
(2) Each player observes its own ships, $\ell^t$, $c^t$, $w^t$, and detects opponent ships w.p.\ $q_{i,v}$. (3) Ships whose endurance is nearly spent are forced toward home\add{ or toward their side's oiler. Ships due for
maintenance or with an engineering casualty are forced home}. (4) Both players choose moves and postures
\emph{simultaneously}. (5) Moves resolve under movement budgets, and endurance updates\add{. Maintenance clocks, engineering
casualties and repairs update, and each oiler moves}. (6) Control indices update,
incidents are drawn, and $\ell$ moves up or down. (7) Stage payoffs are paid, and the game stops if $\ell=L$ (war) or
$t=T=\del{60}\add{180}$.

\begin{table}[H]\centering\scriptsize\setlength{\tabcolsep}{3pt}
\begin{minipage}[t]{0.495\textwidth}\vspace{0pt}
\begin{tabular}{@{}l>{\raggedright\arraybackslash}p{0.44\linewidth}>{\raggedright\arraybackslash}p{0.40\linewidth}@{}}
\toprule
Symbol & Meaning & Value\\\midrule
$T$; $\delta$ & periods (2 days); discount & \del{60}\add{180}; 0.98\\
$N_k$ & fleet size & \del{10 DDG; 8 PLAN, 6 CCG}\add{Table~\ref{tab:fleet}}\\
$D_k$ & endurance (periods) & \del{10 DDG; 8 PLAN; 6 CCG}\add{Table~\ref{tab:fleet}}\\
$F_i$ & movement budget (moves/period) & \del{8 Blue; 10 Red}\add{23 Blue; 40 Red}\\
$w_k$ & strength per hull (rough seas) & \del{1 DDG, PLAN; 0.6 CCG (0.3)}\add{Table~\ref{tab:fleet}}\\
$\alpha$ & assert multiplier & 1.5\\
$\kappa$; $\lambda$ & presence damping; control rate & 1; 0.3\\
$h_0$ & baseline incident hazard & 0.006\\
$h^\B_\As$; $h^\Rd_\As$ & assert hazard & \del{0.03; 0.036/0.012 PLAN/CCG}\add{0.03/0.01; 0.036/0.012 navy/coast guard}\\
\add{$s_h$} & \add{hazard scale} & \add{0.12}\\
\bottomrule
\end{tabular}
\end{minipage}\hfill
\begin{minipage}[t]{0.495\textwidth}\vspace{0pt}
\begin{tabular}{@{}l>{\raggedright\arraybackslash}p{0.44\linewidth}>{\raggedright\arraybackslash}p{0.40\linewidth}@{}}
\toprule
Symbol & Meaning & Value\\\midrule
$L$; $\ell^1$ & war level; starting level & 5; 1\\
$\mu$ & chance quiet period lowers $\ell$ & 0.25 (D3: 0.5 at $\ell\ge3$)\\
$\gamma$; $\eta$ & move; tension cost & 0.1; 1\\
$W$ & war penalty (both sides) & 200\\
$q_{i,v}$ & detection: here / nearby / far & 0.95 / 0.6 / 0.2, $\times0.7$ rough\\
$\Lambda$; $\beta$ & MIP: incident price; persistence & 20; 0.96\\
$\varepsilon_i$ & escalation-risk budget & 0.05 / 0.15 / 0.40\\
$\zeta_D$ & redesign cost per period & D1 0.3; D2 0.5; D3 0\\
\add{$M_{\text{on}}$; $M_{\text{in}}$} & \add{periods operating; in maintenance} & \add{90; 45}\\
\add{$m_0$} & \add{share in maintenance at $t=1$} & \add{1/3}\\
\add{$p_c$; $Z_c$} & \add{casualty chance per period at sea; repair periods} & \add{$1-0.99^2\approx0.02$; 5}\\
\add{$D_o$} & \add{oiler endurance (periods)} & \add{15}\\
\bottomrule
\end{tabular}
\end{minipage}
\caption{Parameter values, all of them our assumptions (no public data). Hazards and $\mu$ were set so that
war is uncommon but not rare (0--34\% of games).}
\label{tab:params}
\end{table}

\subsection{Players, map and time}\label{sec:setup}
\textbf{Players:} $\I=\{\B,\Rd\}$. \del{Blue is the U.S.\ Navy, $\K_\B=\{\text{DDG}\}$. Red is China, with navy ships and
coast guard or maritime militia hulls, $\K_\Rd=\{\text{PLAN},\text{CCG}\}$ (fleets in Table~\ref{tab:params}).}
\add{Blue is the United States with Japan and the Philippines, with eight ship types $\K_\B$. Red is China with North Korea,
with six types $\K_\Rd$ (Table~\ref{tab:fleet}). Each type is navy or coast guard. Coast guard and maritime militia
hulls are grouped as coast guard.}

% !TEX root = updates.tex
\begin{addblock}
\begin{table}[H]\centering\small\setlength{\tabcolsep}{4pt}
\begin{tabular}{@{}llllcccl@{}}
\toprule
Side & Nation & Type $k$ & Ships & $w_k$ & $D_k$ & Real / game & Home ports (real)\\\midrule
Blue & U.S. & US-DDG & destroyers and cruisers, 7th Fleet & 1.0 & 10 & 12 / 4 & YOK 12\\
 & & US-LCS & littoral combat ships$^*$ & 0.5 & 5 & 2 / 1 & GUA 2\\
 & Japan & JP-DDG & JMSDF destroyers & 1.0 & 10 & 26 / 9 & YOK 13, SAS 13\\
 & & JP-FF & JMSDF frigates and escorts & 0.7 & 8 & 16 / 5 & SAS 16\\
 & & JCG & Japan Coast Guard large patrol vessels & 0.6 & 6 & 12 / 4 & ISG 12\\
 & Philippines & PH-FF & Philippine Navy frigates & 0.7 & 8 & 4 / 1 & PHL 4\\
 & & PH-PS & corvettes and offshore patrol vessels$^*$ & 0.4 & 5 & 9 / 3 & PHL 9\\
 & & PCG & Philippine Coast Guard cutters$^*$ & 0.6 & 6 & 6 / 2 & PHL 6\\\midrule
Red & China & CN-DDG & PLAN destroyers & 1.0 & 10 & 30 / 10 & QDO 8, NBO 10, SAN 12\\
 & & CN-FF & PLAN frigates & 0.7 & 8 & 27 / 9 & QDO 7, NBO 9, SAN 11\\
 & & CN-FS & PLAN corvettes$^*$ & 0.4 & 5 & 25 / 8 & QDO 6, NBO 9, SAN 10\\
 & & CCG & China Coast Guard large cutters$^*$ & 0.6 & 6 & 75 / 25 & QDO 19, NBO 26, SAN 30\\
 & North Korea & KP-DDG & destroyers & 1.0 & 8 & 3 / 1 & NAM 3\\
 & & KP-FS & light frigates and corvettes$^*$ & 0.3 & 5 & 11 / 4 & NAM 11\\\midrule
\multicolumn{6}{@{}l}{Totals} & \multicolumn{2}{l}{Blue 87 / 29, Red 171 / 57}\\
\bottomrule
\end{tabular}
\caption{Fleets in theater (no aircraft carriers). Counts are our estimates of the ships each nation keeps in or near
the theater, from open sources \citep{orourke2025,seventhfleet,jmsdf,amti_jcg,phnavy,janes_pcg,kpn}. The Chinese
side is about half of the PLAN and CCG totals, for the share facing this theater. $w_k$ is the strength weight
(destroyers and cruisers highest), $D_k$ the endurance in two-day periods. Types marked $^*$ are small hulls whose
strength is halved in rough seas. ``Game'' is the number of game ships, each standing for three real ships
(\S\ref{sec:hullscale}).}
\label{tab:fleet}
\end{table}

\textbf{Game ships stand for three real ships:}\label{sec:hullscale} the best-response program (\S\ref{sec:mipflow}) takes too
long to solve at full size. We therefore play with one game ship for every
$s=3$ real ships of the same type. A game ship has weight $s\,w_k$, so the local strengths $S_{i,v}$, and with them control and
incident hazards, stay on the scale of the real fleets. Each type's count is divided by $s$ and rounded, with at least one ship
per type, and split over home ports by largest remainder:
\[
N^{\text{game}}_k=\max\!\big(1,\ \mathrm{round}(N_k/s)\big),\qquad
N^{\text{game}}_{k,h}=\lfloor N_{k,h}/s\rfloor+\1\big[h\text{ is among the }N^{\text{game}}_k-\textstyle\sum_{h'}\lfloor N_{k,h'}/s\rfloor\text{ ports with the largest remainders}\big].
\]
Every ship-level rule (endurance, maintenance, engineering casualties, oiler) applies to a game ship as a whole. The
movement budget counts game-ship moves, $F_i=\mathrm{round}(f_i N^{\text{game}}_i)$ with $f_\B=0.8$ and $f_\Rd=0.7$.
\end{addblock}


\textbf{Map:} the theater is an undirected graph $G=(\V,\E)$. Its nodes are home
ports, where ships refuel, and eight contested zones
$\C=\{\text{SEN},\text{MIY},\text{TWN},\text{EOT},\text{LUZ},\text{SCB},\text{STS},\text{PAR}\}$ (Senkakus, Miyako
Strait, Taiwan Strait, east of Taiwan, Luzon Strait, Scarborough Shoal, Second Thomas Shoal, Paracels). Blue's
home ports are $H_\B=\{\text{YOK},\text{GUA}\add{,\text{SAS},\text{ISG},\text{PHL}}\}$ (Yokosuka, Guam\add{, Sasebo, Ishigaki, Subic Bay}) and Red's are $H_\Rd=\{\add{\text{QDO},}\text{NBO},\text{SAN}\add{,\text{NAM}}\}$
(\add{Qingdao, }Ningbo, Sanya\add{, Nampo})\footnote{NBO and SAN stand for the Ningbo--Zhoushan and Hainan (Yulin) basing areas. CCG and militia ports are folded into them.}.
\add{Ports are national. A ship of nation $a$ refuels and does maintenance only at its nation's ports $H_a$: U.S.\ ships
at YOK and GUA, Japanese ships at YOK, SAS and ISG, Philippine ships at PHL, Chinese ships at QDO, NBO and SAN, North
Korean ships at NAM. Design D1 also opens PHL and SAS to U.S.\ ships. The new ports join the zones next to them
(PHL to SCB, LUZ, STS. SAS to SEN, MIY. ISG to SEN, MIY, EOT, TWN. QDO to SEN, NBO, NAM).} An edge joins two nodes a ship can sail between in
one period. The zones form three regions that share weather: North $=\{\text{SEN},\text{MIY}\}$, Center
$=\{\text{TWN},\text{EOT},\text{LUZ}\}$, South $=\{\text{SCB},\text{STS},\text{PAR}\}$. Zone $v$ is worth $V^i_v$
to player $i$ (Table~\ref{tab:zones}, e.g.\ Scarborough Shoal and the Paracels matter more to Red).

\textbf{Time:} periods $t=1,\dots,T$ of two days, $T=$ \del{60 (about four months)}\add{180 (one year)}, discounted by $\delta$ per period. All ships start in port with full endurance, \del{split evenly between each side's two ports}\add{at their home ports (Table~\ref{tab:fleet})}\add{. A third of each fleet starts in maintenance (\S\ref{sec:readiness})}.

\begin{table}[H]\centering\small
\begin{minipage}[t]{0.56\textwidth}\vspace{0pt}\centering
\begin{tabular}{@{}lcccccccc@{}}
\toprule
Zone & TWN & SEN & MIY & EOT & LUZ & SCB & STS & PAR\\\midrule
$V^\B_v$ & 10 & 6 & 4 & 4 & 6 & 5 & 5 & 2\\
$V^\Rd_v$ & 10 & 6 & 3 & 4 & 5 & 7 & 6 & 6\\
$c^1_v$ & 0 & 0.2 & 0.3 & 0.2 & 0.1 & $-0.2$ & 0 & $-0.4$\\
\bottomrule
\end{tabular}
\end{minipage}\hfill
\begin{minipage}[t]{0.42\textwidth}\vspace{0pt}\centering
\begin{tabular}{@{}lcccc@{}}
\toprule
$\omega$ & calm & nominal & typhoon & extreme\\\midrule
$\pi_\omega$ & 0.25 & 0.40 & 0.25 & 0.10\\
$p_s(\omega)$ & 0 & 0.02 & 0.06 & 0.12\\
$p_r(\omega)$ & 0 & 0.20 & 0.35 & 0.50\\
\bottomrule
\end{tabular}
\end{minipage}
\caption{Our assumed values. Left: zone values and starting control. Right: scenario weight $\pi_\omega$ and per-period chances of a storm
starting $p_s(\omega)$ and of rough seas $p_r(\omega)$.}
\label{tab:zones}
\end{table}

\subsection{State, information and weather}\label{sec:state}
\textbf{State:} $s^t=(x^t,d^t,c^t,\ell^t,w^t\add{,r^t,z^t,e^t,o^t,d^{o,t}})$. Here $x^t_{k,v}\in\Z_{\ge0}$ is the number of type-$k$ ships at node
$v$. The \emph{endurance} $d^t_j\le D_k$ is the number of periods ship $j$ can still stay at sea (down one per period
at sea, refilled after a full period in port). The \emph{control index} $c^t_v\in[-1,1]$ of each zone is $+1$ when
Blue is fully in charge and $-1$ when Red is. The \emph{escalation level} $\ell^t\in\{0,\dots,L\}$ runs from calm to
armed conflict. The sea state of region $r$ is $w^t_r\in\{\text{calm},\text{rough},\text{storm}\}$. \add{Each ship $j$ also carries
a readiness state: $r^t_j$ is the number of operating periods left before it is due for maintenance, $z^t_j$ is the
number of periods it still has to spend in port (maintenance or repair, $0$ when it is available), and
$e^t_j\in\{0,1\}$ is $1$ if it has an engineering casualty and is heading home. Each player's oiler is at node
$o^t_i$ with $d^{o,t}_i$ periods of its own fuel left (\S\ref{sec:readiness}).}

\textbf{Information:} each player sees its own ships and the public $c^t$, $\ell^t$, $w^t$, but not where the
opponent's ships are now. Each period it detects the opponent at node $v$ with probability $q_{i,v}$
(Table~\ref{tab:params}, for an own ship there, a ship next door or $v$ near a home port, or satellite only), seeing the count
of each type and the posture. Its \emph{belief} $\hat x^t_{-i,v}$ is the last count it detected at $v$ (exact at
the start, when all ships are in port), so moves are chosen at information sets that do not reveal where the opponent is now.

\textbf{Weather:} a random factor neither side controls. In each of four scenarios (not calendar seasons)
$\omega\in\Omega=\{$calm, nominal, typhoon, extreme$\}$ with weights $\pi_\omega$, each period a region
not already in a storm enters one with probability $p_s(\omega)$, and is otherwise rough with probability
$p_r(\omega)$ (Table~\ref{tab:zones}). A storm lasts three periods. It closes the region's zones (no entry, ships
inside stay), stops incidents there and freezes their control. Rough seas halve \del{CCG}\add{small-hull} strength and degrade detection.

\subsection{Actions and constraints}\label{sec:actions}
\textbf{Actions:} each period both players choose, simultaneously, a destination for each ship (stay, or sail one
edge) and a posture $p_{i,v}\in\{\Pt,\As\}$ in each zone $v$ where they will have ships. \emph{Patrol} is routine
presence. \emph{Assert} stands for the \S1 actions, treated as one posture. Asserting multiplies the side's strength
there by $\alpha$ but makes incidents more likely (\S\ref{sec:dynamics}).

\textbf{Constraints:} with $f^t_{k,uv}$ the number of type-$k$ ships sailing $u\to v$ in period $t$ and $\mathcal N(u)$ the neighbors of $u$,
\begin{align*}
&x^{t+1}_{k,u}=x^t_{k,u}+\textstyle\sum_{v\in\mathcal N(u)}f^t_{k,vu}-\sum_{v\in\mathcal N(u)}f^t_{k,uv}
&&\text{(flow conservation: no ships are added or removed)}\\
&\textstyle\sum_{k\in\K_i}\sum_{(u,v)}f^t_{k,uv}\le F_i
&&\text{(movement budget: at most $F_i$ moves per period, Table~\ref{tab:params})}
\end{align*}
where the theater-wide movement budget $F_i$ stands in for operational requirements and station-keeping.
\emph{Endurance:} a ship may not sail to a node from which it could not \del{get home}\add{reach a refuel point} on its remaining endurance,
unless the move brings it closer to \del{home}\add{one}. \add{The refuel points of player $i$ are its home ports and its oiler when the oiler is in a zone,
$\mathcal R^t_i=H_i\cup\{o^t_i : o^t_i\in\C\}$. A ship uses whichever is closer, at distance
$\mathrm{dist}(v,\mathcal R^t_i)=\min_{u\in\mathcal R^t_i}\mathrm{dist}_i(v,u)$.} A ship whose endurance has fallen to its
distance from \del{home}\add{the nearest refuel point} is on
\emph{forced return}. It sails one edge toward \del{port}\add{that point} each period, then spends one full period in port \add{or alongside the oiler} to refit,
after which $d_j$ resets to $D_k$. Forced moves use the movement budget first, and the player decides only for its
\emph{free} ships. \emph{Ports, storms and infeasible orders:} no ship may enter opponent home ports or storm
regions, and orders over budget are dropped.

% !TEX root = updates.tex
\begin{addblock}
\subsection{Readiness: maintenance, engineering casualties and oilers}\label{sec:readiness}
A year-long game needs the things that keep a real fleet from being at sea all the time.

\textbf{Maintenance:} a ship operates $M_{\text{on}}=90$ periods (six months) and then spends $M_{\text{in}}=45$
periods (three months) in maintenance at a home port, so in steady state a third of each fleet is unavailable.
While a ship is available its maintenance clock runs down, $r^{t+1}_j=r^t_j-1$. When $r^t_j\le0$ the ship is on
forced return to its nearest home port (not to the oiler). On arrival it enters maintenance, $z_j=M_{\text{in}}$, and
$z_j$ falls by one each period. When $z_j$ reaches $0$ the ship is available again with $r_j=M_{\text{on}}$ and full
endurance $d_j=D_k$. A ship with $z_j>0$ stays in port and is not part of either player's decision.

\textbf{Staggered start:} at $t=1$ a share $m_0=1/3$ of each player's ships, chosen at random, is in maintenance with
$z^1_j$ drawn uniformly from $\{1,\dots,M_{\text{in}}\}$. Every other ship has $r^1_j$ drawn uniformly from
$\{1,\dots,M_{\text{on}}\}$, so ships come due at different times and not all at once. These draws use their own
random stream, so a given seed has the same maintenance schedule for every pair of policies.

\textbf{Engineering casualties:} each day at sea a ship suffers an engineering casualty with probability $1\%$, so
in a two-day period the chance is $p_c=1-0.99^2\approx0.02$. Then $e_j=1$ and the ship is on forced return to its
nearest home port. On arrival it is repaired for $Z_c=5$ periods (ten days), $z_j=Z_c$, after which $e_j=0$ and
$d_j=D_k$. A ship that has a casualty while heading in for maintenance is repaired during the maintenance period.

\textbf{Oilers:} each player has one oiler. It is not commanded ship by ship. It follows a fixed rule that keeps it
where its side is doing well. The side's \emph{best region} is the one with the highest value-weighted control from
that side's point of view,
\[
r^{*,t}_i=\arg\max_{r\in\{\text{N},\text{C},\text{S}\}}\ \sigma_i\,\frac{\sum_{v\in r}V^i_v\,c^t_v}{\sum_{v\in r}V^i_v}.
\]
The oiler picks a zone of $r^{*,t}_i$ uniformly at random as its target, sails one edge per period toward it, and
picks a new random target when it arrives or when the best region changes. So it wanders among the zones of the
best region and follows when that region changes. Its own fuel $d^{o}_i$ falls by one per period away from port.
When $d^{o,t}_i\le\mathrm{dist}_i(o^t_i,H_i)$ it heads for its nearest home port, and one period there refills it to
$D_o=15$ (30 days). A ship that spends a full period in the oiler's zone without moving is refueled to $D_k$, as
in port. The oiler obeys the same storm and port rules as warships. It adds no strength and cannot be attacked yet.
\end{addblock}

\subsection{How the state changes: control, incidents and escalation}\label{sec:dynamics}
After the moves, steps 1--3 apply in each contested zone $v$ not in a storm. Step 4 updates $\ell$ once for the theater.

\textbf{1. Local strength:} each side's weight of presence in the zone is
$S_{i,v}=\alpha^{\1[p_{i,v}=\As]}\sum_{k\in\K_i}w_k\,x_{k,v}$, its hulls weighted by $w_k$ (a CCG hull counts less
than a navy ship, less still in rough seas), times $\alpha$ when asserting.

\textbf{2. Control:} the \emph{local balance}
$\phi_v=(S_{\B,v}-S_{\Rd,v})/(S_{\B,v}+S_{\Rd,v}+\kappa)$ lies between $-1$ and $1$ and is positive when Blue is
locally stronger. The constant $\kappa=1$ keeps a small presence from counting as full control (e.g. one DDG alone gives
$\phi_v=0.5$). Control is earned by staying present and fades when a side leaves. Each period it moves a fraction
$\lambda=0.3$ of the way toward the local balance,
\begin{equation}
c^{t+1}_v=(1-\lambda)\,c^t_v+\lambda\,\phi_v .\label{eq:control}
\end{equation}

\textbf{3. Incidents:} an incident (a collision, an unsafe interaction, negative media coverage) can happen only
where both sides have ships. It occurs with probability $\rho_v=1-e^{-h_v}$, where the \emph{hazard} $h_v$ is
\begin{equation}
h_v=\add{s_h}\big(h_0+h^\B_{\As\add{,v}}\,\1[p_{\B,v}=\As]+h^\Rd_{\As,v}\,\1[p_{\Rd,v}=\As]\big)\cdot\min\{S_{\B,v},S_{\Rd,v}\}
\cdot\big(1+\ell^t/2\big).\label{eq:hazard}
\end{equation}
(i) A small baseline hazard $h_0$, plus a term for each side that asserts. \del{Red's assert
hazard $h^\Rd_{\As,v}=\theta_v h^\Rd_{\As,\text{PLAN}}+(1-\theta_v)h^\Rd_{\As,\text{CCG}}$ blends its navy and coast
guard values by $\theta_v$, the PLAN share of Red's strength at $v$}\add{Each side's assert hazard
$h^i_{\As,v}=\theta_{i,v}h^i_{\As,\text{navy}}+(1-\theta_{i,v})h^i_{\As,\text{cg}}$ blends its navy and coast guard values by
$\theta_{i,v}$, the navy share of its strength at $v$}, so coast guard pressure is less escalatory.
\add{The scale $s_h$ multiplies all hazards. It is set so that rule-based play on both sides ends in war in about 5\% of
one-year games in nominal weather.}
(ii) More ships in contact mean more encounters, limited by the smaller force. (iii) Incidents are likelier when
tension is high ($2.5\times$ at $\ell=3$).

\textbf{4. Escalation:} each incident raises $\ell$ by one. A period with no incident anywhere lowers $\ell$ by one
with probability $\mu$ (Table~\ref{tab:params}). If $\ell$ reaches $L$, the crisis has become an armed conflict and the game ends.

\subsection{Payoffs}\label{sec:payoffs}
\textbf{Stage payoff:} at the end of period $t$, player $i$ receives
\begin{equation}
g_i^t=\sigma_i\textstyle\sum_{v\in\C}V^i_v\,c^{t+1}_v-\gamma\,m^t_i-\eta\,\ell^{t+1}-\zeta_D\,\1[i=\B].\label{eq:stage}
\end{equation}
\emph{Control:} $\sigma_\B=+1$ and $\sigma_\Rd=-1$ turn the shared control index into each side's gain (e.g.
Blue's control term is $V^\B_v c_v$, positive where Blue holds the zone and negative where Red does).
\emph{Fuel:} $m^t_i$ ship moves in period $t$, at $\gamma$ each.
\emph{Tension:} $\eta$ per escalation level per period, for diplomatic and economic strain.
\emph{Redesign cost:} $\zeta_D$ is Blue's redesign cost per period (\S\ref{sec:designs}).

\textbf{Total payoff:} the game ends at $\tau$, the first period with $\ell^\tau=L$ (war), or at $T$. Player $i$'s
payoff is the discounted sum of its stage payoffs, minus a war penalty $W$ (Table~\ref{tab:params}) if the game ends in war:
\begin{equation}
u_i=\sum_{t=1}^{\tau}\delta^{t-1}g_i^t-W\,\1[\ell^\tau=L].\label{eq:payoff}
\end{equation}
Both sides pay $W$, so the game is general-sum, not zero-sum. Both prefer to avoid war. For scale,
Blue earns at most $\sum_vV^\B_v=42$ in a period in which it fully controls every zone, so a war ($W=200$) costs roughly five such periods.

\subsection{How a player decides: the best-response MIP with a risk budget}\label{sec:mip}
Our main decision rule, \emph{MIP-BR}, is a best response to the opponent's believed positions that caps incident
risk, written as a mixed-integer program and solved with Gurobi each period. It takes as given its beliefs
$\hat x_{-i}$, the sea state, $\ell^t$, its own ships, and, in each zone, the opponent posture it saw on its last
detection there (patrol if never detected). Forced-return ships are fixed ($\bar x_{k,v}$, by destination).

\textbf{Decision variables:} $y_{j,w}\in\{0,1\}$ is 1 if free ship $j$ goes to node $w$ (only moves allowed by
\S\ref{sec:actions}). Control and incident risk are nonlinear in ship counts, so we list every count a zone could
end up with. Option $b_{v,(n,p)}\in\{0,1\}$ is 1 if the player will have $n=(n_k)_{k\in\K_i}$ hulls in zone $v$ with posture
$p$. Each option's value and risk are computed in advance, which keeps the program linear.

\textbf{Value of an option:}\vspace{-4pt}
\[
a_v(n,p)=V^i_v\,\beta\,\sigma_i\phi_v(n,p)-\Lambda\,\rho_v(n,p),\qquad\beta=\frac{\lambda}{1-\delta(1-\lambda)}\approx0.96 .
\]
with $\phi_v,\rho_v$ (\S\ref{sec:dynamics}) at believed opponent positions. Here $\beta$ is the discounted effect
on control of one period of balance, and $\Lambda$ is the incident price.

\textbf{The program:}
\begin{align*}
\max_{y,\,b}\ \ &\textstyle\sum_{v\in\C}\sum_{(n,p)}a_v(n,p)\,b_{v,(n,p)}
-\gamma\sum_{j}\sum_{w\ne\text{loc}(j)}y_{j,w}-1000\,\xi
&&\text{(control value $-$ incident cost $-$ fuel cost)}\\
\text{s.t.}\ \ &\textstyle\sum_w y_{j,w}=1\quad\forall j
&&\text{(each free ship gets one destination)}\\
&\textstyle\sum_{j}\sum_{w\ne\text{loc}(j)}y_{j,w}\le F_i-\text{(forced moves)}
&&\text{(movement budget)}\\
&\textstyle\sum_{(n,p)}b_{v,(n,p)}=1\quad\forall v
&&\text{(one option per zone)}\\
&\textstyle\sum_{(n,p)}n_k\,b_{v,(n,p)}=\bar x_{k,v}+\sum_{j\text{ of type }k}y_{j,v}\quad\forall v,k
&&\text{(the option matches the moves)}\\
&\textstyle\sum_{v\in\C}\sum_{(n,p)}\rho_v(n,p)\,b_{v,(n,p)}\le\varepsilon_i\,(1-\ell^t/L)+\xi
&&\text{(escalation-risk budget)}\\
&y,\,b\ \text{binary},\ \xi\ge0.
\end{align*}
\textbf{The escalation-risk budget:} the left side of the last constraint is the expected number of incidents this
period. The cap $\varepsilon_i$ shrinks to zero as $\ell$ approaches war and defines the feasible region, the
allocations the player considers safe enough. The proposal's conservative, moderate and aggressive commanders become
three budgets (Table~\ref{tab:params}), and RQ1 varies $\varepsilon_i$. A ship on forced return may have to pass through
a zone the opponent occupies. That risk cannot be avoided, and near war it can exceed the cap. The slack $\xi\ge0$,
penalized at 1000 per unit, lets the program exceed the cap by only that amount, so it always has a solution.

% !TEX root = updates.tex
\begin{addblock}
\subsection{The best-response MIP for the Workshop 2 fleets}\label{sec:mipflow}
The program above is exact, but it does not scale to the Workshop 2 fleets. This subsection shows why and gives the
version we use. The rules of the game (\S\ref{sec:actions}--\S\ref{sec:payoffs}) do not change. Only the way the
player's decision problem is written down changes, and it stays exact.

\textbf{Why the Workshop 1 program breaks:} it lists every count vector a zone could end up with. If $A_{k,v}$ ships
of type $k$ could arrive at zone $v$ this period, the number of options in that zone is
\[
|\{b_{v,(n,p)}\}| \;=\; 2\prod_{k\in\K_i}\big(A_{k,v}+1\big).
\]
In Workshop 1, $|\K_\B|=1$ and $|\K_\Rd|=2$, so this was a few hundred. With the eight Blue types at full size and, for example,
$A_{k,v}=(12,13,16,12,2,4,9,6)$, it is
$2\cdot13\cdot14\cdot17\cdot13\cdot3\cdot5\cdot10\cdot7\approx8\times10^{7}$
options in one zone. The model cannot be built, whatever the license or the run time.

\textbf{Step 1, moves as flows:} free ships of the same type, at the same node, with the same set of allowed
destinations are interchangeable. Call such a set a group $g$, with type $k(g)$, location $u(g)$, allowed
destinations $\mathcal W_g$ and size $|g|$. Replace the per-ship binaries $y_{j,w}$ by integer flows
\[
f_{g,w}\in\Z_{\ge0}\ \ (w\in\mathcal W_g),\qquad \textstyle\sum_{w\in\mathcal W_g}f_{g,w}=|g| .
\]
These are the flows $f^t_{k,uv}$ of \S\ref{sec:actions}. Every assignment of ships gives a flow and every flow gives an
assignment, so this step changes nothing.

\textbf{Step 2, options by navy and coast guard strength level:} a zone's value depends on the player's ships there only
through two numbers, its navy strength and its coast guard strength. Control uses their sum. The assert hazard also
uses the navy share, because it blends the navy and coast guard values,
$h^i_{\As,v}=\theta_{i,v}h^i_{\As,\text{navy}}+(1-\theta_{i,v})h^i_{\As,\text{cg}}$ (\S\ref{sec:dynamics}).
Every weight $s\,w_k$, halved or not in rough seas, is a whole multiple of the unit $u=0.05\,s$ ($u=0.15$ for $s=3$), so
$u_k(w_r)=s\,w_k(w_r)/u$ is an integer. The navy and coast guard levels in zone $v$ are
\[
n_v=\sum_{k\ \text{navy}}u_k\Big(\bar x_{k,v}+\sum_{g:\,k(g)=k}f_{g,v}\Big),\qquad
c_v=\sum_{k\ \text{coast guard}}u_k\Big(\bar x_{k,v}+\sum_{g:\,k(g)=k}f_{g,v}\Big),
\]
so the own un-asserted strength is $R_v=u\,(n_v+c_v)$ and the navy share is $\theta_{i,v}=n_v/(n_v+c_v)$. The option
$b_{v,(n,c,p)}\in\{0,1\}$ means ``navy level $n$, coast guard level $c$ and posture $p$ in zone $v$''. Only levels that the
fixed ships and the arriving groups can actually produce are listed. The option values are the same formulas as before,
\[
a_v(n,c,p)=V^i_v\,\beta\,\sigma_i\phi_v(n,c,p)-\Lambda\rho_v(n,c,p),\qquad S_{i,v}=\alpha^{\1[p=\As]}\,u\,(n+c),
\]
with the hazard in $\rho_v$ using $\theta_{i,v}=n/(n+c)$. Both control and incident risk are therefore exactly those of the
Workshop 1 program. A patrol option does not depend on the split between $n$ and $c$, because the navy share enters only the
player's own assert hazard. So patrol options need only the total level, which keeps the program small.

\textbf{The program:}
\begin{align*}
\max_{f,\,b}\ \ &\textstyle\sum_{v\in\C}\sum_{(n,c,p)}a_v(n,c,p)\,b_{v,(n,c,p)}
-\gamma\sum_{g}\sum_{w\ne u(g)}f_{g,w}-1000\,\xi
&&\text{(control value $-$ incident cost $-$ fuel cost)}\\
\text{s.t.}\ \ &\textstyle\sum_{w\in\mathcal W_g}f_{g,w}=|g|\quad\forall g
&&\text{(every free ship gets one destination)}\\
&\textstyle\sum_{g}\sum_{w\ne u(g)}f_{g,w}\le F_i-\text{(forced moves)}
&&\text{(movement budget)}\\
&\textstyle\sum_{(n,c,p)}b_{v,(n,c,p)}=1\quad\forall v
&&\text{(one option per zone)}\\
&\textstyle\sum_{(n,c,p)}n\,b_{v,(n,c,p)}=n_v,\quad \sum_{(n,c,p)}c\,b_{v,(n,c,p)}=c_v\quad\forall v
&&\text{(the option matches the moves)}\\
&\textstyle\sum_{v\in\C}\sum_{(n,c,p)}\rho_v(n,c,p)\,b_{v,(n,c,p)}\le\varepsilon_i\,(1-\ell^t/L)+\xi
&&\text{(escalation-risk budget)}\\
&f\ \text{integer},\ b\ \text{binary},\ \xi\ge0.
\end{align*}
Compared with \S\ref{sec:mip}, $y$ becomes $f$, and the index $n$ (a count for each type) becomes $(n,c)$ (a navy and a
coast guard strength level). Everything else is unchanged. In the code, the options of one posture (and, for assert, one coast guard
level) are entered as a piecewise-linear function of the level with a breakpoint at every level that can occur, which is
the same program with fewer variables.

\textbf{Check:} at every solve we checked, the program's objective equals the value of its plan recomputed by the
game engine (\texttt{plan\_value}), and it is never below the value of the greedy plan for the same state.

\textbf{Run time and the game-ship scale:} at full size, even an earlier version of this program that approximated the
navy share took 13 seconds to 4 minutes per solve, and a one-year game needs 360 solves (two players, 180 periods).
With each game ship standing for three real ships (\S\ref{sec:hullscale}), a year of MIP-BR against MIP-BR takes 30 to 39
minutes on one core. The experiments (about 12{,}400 games) therefore run as independent jobs on MIT's Engaging cluster.
\end{addblock}


\subsection{The template game and the policies we compare}\label{sec:policies}
\textbf{Template game:} the full game has far too many strategies to solve for an equilibrium directly, so we
also study a small strategic-form version. Each player picks one of 14 \emph{templates}, which pair one of seven ways to split
\del{70\%}\add{two thirds (the share out of maintenance)} of its fleet across North/Center/South (60/20/20, 40/40/20 or 34/33/33 percent and their permutations) with one
posture (P or A) in all its zones. Ships are spread evenly within a region, Red's PLAN and CCG in proportion, and
\del{fractional hulls are allowed}\add{only whole ships are used. With $N^{\text{dep}}_k$ the deployable ships of type $k$
rounded to a whole number and $\pi_r$ the template's share for region $r$, the region counts are the
largest-remainder rounding of $N^{\text{dep}}_k\pi_r$. Each region gets $\lfloor N^{\text{dep}}_k\pi_r\rfloor$ ships,
and the ships left over go one at a time to the regions with the largest fractional parts. The same rule splits a
region's ships over its zones}. If both kept their
templates, control would settle at $c_v=\phi_v$ (set $c^{t+1}_v=c^t_v$ in \eqref{eq:control}), so a template
pair's payoff is that settled control minus the incident cost, in calm seas at $\ell=1$:
$M_i(a_\B,a_\Rd)=\sigma_i\textstyle\sum_{v\in\C}V^i_v\,\phi_v-\Lambda\sum_{v\in\C}\rho_v$. This gives a finite game $\Gamma=\langle\I,(S_i),(M_i)\rangle$ with $|S_i|=14$, which we analyze with iterated
elimination of strictly dominated strategies (IESDS), pure-NE enumeration, and the maximin and correlated-equilibrium LPs.

\textbf{Policies} (\S4): \emph{Random} picks moves and postures at random, as a benchmark. \emph{Rule-based}
(moderate commander) sends ships where value per unit of opponent strength is highest, asserts only where it believes
it outnumbers the opponent and $\ell\le2$, then pulls back until its expected incidents fit $\varepsilon=0.15\,(1-\ell/L)$.
\emph{Greedy} is a fast heuristic for the MIP (same objective and budget, no solver). \emph{MIP-BR}$(\varepsilon)$
solves the MIP exactly. \emph{Equilibrium} plays its part of the pure NE of $\Gamma$ with the higher total payoff
(our choice), steering free ships toward the template's split (more ships to a region's more valuable zones).
\emph{Maximin} draws a template every three periods from its maximin mixed strategy of $\Gamma$. The two template policies switch to patrol when $\ell>3$ and drop
assertions over their risk budget.

\subsection{Resilient redesigns (Blue)}\label{sec:designs}
A redesign changes Blue's system. D1 is structural (an alternative
path), D2 and D3 operational (monitoring, rapid response). Each is compared with the baseline D0 in every scenario.
\textbf{D1 dispersed basing:} a third Blue port at Subic Bay, linked to SCB, LUZ and STS, giving
alternative paths and shorter trips South. Blue starts 4/3/3 at YOK/GUA/Subic. \textbf{D2 redundant sensing:} Blue
detects in every zone with probability at least $0.6$, even in bad weather. \textbf{D3 de-escalation measures:} once $\ell\ge3$, a quiet period lowers $\ell$ with probability $0.5$ instead of $0.25$, and assert
hazards are halved, which helps Red as much as Blue. Costs $\zeta_D$ are in Table~\ref{tab:params}.

\textbf{Simplifying assumptions:} there are two unitary players, control is a force-ratio index, incidents are the
only path to war, beliefs are last-seen counts, and best responses look one period ahead. \add{Oilers follow a fixed
rule and cannot be targeted, maintenance takes place only at home ports, and engineering casualties are independent
from ship to ship.}

%======================================================================
\section{Preliminary results}
%======================================================================
\add{\textbf{Note:} every number in this section is new (one-year game, Workshop 2 fleets at one game ship per three
real ships, readiness, exact MIP, $s_h=0.12$). Workshop 1 values are shown in gray for comparison. Payoffs are discounted
sums over 180 periods instead of 60, so they are larger in size than in Workshop 1.}

% !TEX root = report.tex
\S4.1 is analytical. The rest use 200 seeds per configuration.
Unless stated, Red plays MIP-BR with $\varepsilon_\Rd=0.15$ in nominal weather.

\subsection{Template game and a small example: patrol or assert, and sustainable restraint}
IESDS reduces $\Gamma$ from $14\times14$ to $2\times4$. Blue keeps only the center-heavy split (20/60/20), and Red
keeps the south-heavy and center/south-heavy splits. Of the two pure NE, the \emph{equilibrium} policy plays (Blue center-heavy
assert, Red south-heavy patrol), $(0.39,-9.27)$, over (Blue patrol, Red center/south assert), $(-14.44,2.71)$.
Without the escalation cost, the zero-sum maximin LP has a pure equilibrium where \emph{both} sides concentrate on
the Center and assert. The shared incident cost is what makes at most one side assert.
Fixing Blue center-heavy and Red south-heavy, each side only chooses P or A. The resulting $2\times2$ sub-matrix of $\Gamma$, normalized so (P,P) $=(0,0)$, has
(P,A) $=(-9.61,\,1.20)$, (A,P) $=(0.28,\,-10.45)$ and (A,A) $=(-11.41,\,-11.33)$ (Blue, Red).
Asserting against a patroller pays a little (one-period gain $G_\B=0.28$, $G_\Rd=1.20$), but mutual assertion costs more than
patrolling alone, so $BR_i(\Pt)=\As$ and $BR_i(\As)=\Pt$. There are two pure NE, (A,P) and (P,A), and a mixed NE (Blue asserts w.p.\
$0.58$, Red $0.13$) with mutual assertion w.p.\ $0.08$ and payoffs $v_\B=-1.30$, $v_\Rd=-6.02$. The welfare-maximizing correlated equilibrium (CE LP)
puts \emph{zero} on (A,A). A mediator (e.g.\ a hotline) removes mutual assertion.

In the repeated game, suppose both patrol, and any assertion makes both fall back to the
mixed NE for good. By the one-stage deviation principle, patrolling is optimal iff
$(1-\delta)G_i+\delta v_i\le0$, i.e.\ $\delta\ge G_i/(G_i-v_i)$, which is $0.18$ for Blue and $0.17$ for Red. With
$\delta=0.98$, \textbf{mutual restraint is a subgame perfect equilibrium}, though our one-period-ahead policies
cannot see this.

\subsection{RQ1: control versus war risk, and the risk-budget game}
\noindent\begin{minipage}[t]{0.40\textwidth}\vspace{0pt}\centering\footnotesize\setlength{\tabcolsep}{2pt}
\begin{tabular}{l|ccc}
\toprule
$\B\backslash\Rd$ & cons.\ $0.05$ & mod.\ $0.15$ & aggr.\ $0.40$\\\midrule
cons.\ $0.05$ & $-194,\ 199$ & $-303,\ 315$ & $-308,\ 318$\\
              & war $0.10$ & $0.11$ & $0.12$\\
mod.\ $0.15$  & $18,\ -36$ & $-144,\ 101$ & $-192,\ 140$\\
              & war $0.15$ & $0.20$ & $0.21$\\
aggr.\ $0.40$ & $74,\ -80$ & $-83,\ 16$ & $\mathbf{-151,\ 39}$\\
              & war $0.14$ & $0.25$ & $\mathbf{0.34}$\\
\bottomrule
\end{tabular}
\captionof{table}{The risk-budget game: $(u_\B,u_\Rd)$ and war rate when both sides choose $\varepsilon$
(s.e.\ $\approx 5$--$12$). Bold: the unique NE.}\label{fig:rq1}
\end{minipage}\hfill
\begin{minipage}[t]{0.58\textwidth}\vspace{0pt}
\textbf{Frontier:} Raising $\varepsilon_\B$ from $0.02$ to $0.15$ lifts Blue's control share from $-0.25$ to
$-0.05$ while the war rate rises from $0.07$ to $0.20$ (Red fixed at $0.15$). Going to $0.40$ adds only $0.06$ of control. Most of what a looser budget buys is bought by $\varepsilon\approx0.15$--$0.25$.

\smallskip\textbf{Risk-budget game} (Table~\ref{fig:rq1})\textbf{:} each side picks a budget. Whatever the other
side picks, ``aggressive'' pays each side the most, so both pick it, and
(aggressive, aggressive) is the unique pure NE. It also has the \emph{highest} war rate ($0.34$), and both sides
would do at least as well if both chose moderate ($-144,\,101$ vs.\ $-151,\,39$, war $0.20$). This is a prisoner's
dilemma, so restraint must come from repetition or coordination (\S4.1).
\end{minipage}

\subsection{RQ2: strategic versus heuristic play}
\begin{table}[!ht]\centering\footnotesize\setlength{\tabcolsep}{4pt}\renewcommand{\arraystretch}{0.9}
\begin{tabular}{l|ccc}
\toprule
Blue $\backslash$ Red & rule-based (moderate) & MIP-BR & equilibrium\\\midrule
random & $-283$ / $-0.14$ / $0.15$ & $-319$ / $-0.16$ / $0.33$ & $-164$ / $-0.08$ / $0.06$\\
rule-based (moderate) & $-122$ / $-0.05$ / $0.04$ & $-184$ / $-0.09$ / $0.06$ & $35$ / $+0.06$ / $0.04$\\
greedy & $-70$ / $+0.01$ / $0.12$ & $-137$ / $-0.05$ / $0.14$ & $65$ / $+0.08$ / $0.03$\\
MIP-BR & $-75$ / $+0.01$ / $0.13$ & $-144$ / $-0.05$ / $0.20$ & $53$ / $+0.08$ / $0.08$\\
equilibrium & $-72$ / $-0.01$ / $0.07$ & $\mathbf{-116}$ / $-0.03$ / $0.10$ & $21$ / $+0.05$ / $0.06$\\
maximin & $-170$ / $-0.09$ / $0.04$ & $-221$ / $-0.12$ / $0.08$ & $-35$ / $-0.01$ / $0.00$\\
\bottomrule
\end{tabular}
\caption{Blue payoff $u_\B$ / control share / war rate (s.e.\ of $u_\B$ is 1--7). Nominal weather, design D0. The
MIP-BR column is each policy's worst case.}
\label{tab:policies}
\end{table}
Greedy, MIP-BR and equilibrium beat the random and rule-based benchmarks (Table~\ref{tab:policies}). The
\textbf{equilibrium} policy has the best worst case ($-116$) and, against MIP-BR, half the war rate of MIP-BR itself
($0.10$ vs.\ $0.20$). Greedy and MIP-BR do best against weaker opponents but escalate more. \textbf{Maximin} is the
safest (war rate $\le0.08$) but gives up the most control, guarding against templates Red rarely plays.
Greedy matched the MIP optimum in 67\% of 1{,}524 Blue decisions but, unlike the MIP, broke the risk budget in 14\%.

\subsection{RQ3: robustness to weather and resilient designs}
% !TEX root = updates.tex
\noindent\begin{minipage}[t]{0.47\textwidth}\vspace{0pt}\centering\footnotesize\setlength{\tabcolsep}{2pt}
{\begin{tabular}{l|cccc|cc}
\toprule
Design & calm & nom. & typh. & extr. & exp. & worst\\\midrule
D0 baseline   & $-458$/.34 & $-430$/.26 & $-399$/.21 & $-366$/.16 & $-423$ & $-458$\\
D1 dispersed  & $-474$/.34 & $-421$/.23 & $-394$/.19 & $-368$/.11 & $-422$ & $-474$\\
D2 sensing    & $-479$/.34 & $-447$/.28 & $-424$/.28 & $-386$/.12 & $-443$ & $-479$\\
D3 de-escalation & $\mathbf{-444}$/.10 & $\mathbf{-395}$/.11 & $\mathbf{-382}$/.06 & $\mathbf{-350}$/.04 & $\mathbf{-400}$ & $\mathbf{-444}$\\
\bottomrule
\end{tabular}}
\captionof{table}{Blue designs by weather scenario (MIP-BR vs.\ MIP-BR): $u_\B$ / war rate (s.e.\ $\approx6$--$9$).}\label{fig:rq3}
\smallskip\normalsize\raggedright We score each policy or design $x$ by its expected payoff
$\sum_\omega\pi_\omega J(x,\omega)$ and worst case $\min_\omega J(x,\omega)$.
\end{minipage}\hfill
\begin{minipage}[t]{0.51\textwidth}\vspace{0pt}
\textbf{Weather:} as in Workshop 1, bad weather helps Blue. Every policy except MIP-BR with $\varepsilon=0.40$ does
worst in \emph{calm} seas, since rough seas halve small-hull strength, which is about half of Red's strength but less than a quarter of
Blue's, and storms
pause operations and incidents (MIP-BR war rate $0.34$ calm, $0.16$ extreme). The top of the ranking is the same in all
four scenarios: MIP-BR with $\varepsilon=0.40$, then $\varepsilon=0.15$, then greedy. Only the three weakest policies
(rule-based, equilibrium, maximin) change places. The best policy on average ($\varepsilon=0.40$, expected $3$) is also
best in its worst case ($-88$, extreme seas).
\end{minipage}

\textbf{Designs:} de-escalation measures (D3) are again best in expected and worst-case terms. They cut the war rate
from $0.26$ to $0.11$ (nominal) \wone{$0.20\to0.05$} and raise Blue's payoff by $35$. Red's payoff rises too
($385\to417$), since fewer wars help both sides. We still charge them nothing, so part of this win is by assumption.
Dispersed basing (D1) leaves Blue's payoff where it was ($-421$ against $-430$) and no longer cuts Red's
\wone{$101\to16$}, now $385\to371$. Redundant sensing (D2) is worse than the baseline ($-447$), paying its cost for no
measurable benefit.
\textbf{Limitations:} parameters are assumptions, policies look one period ahead, and each game ship stands for three
real ships. Payoff differences below about $20$ are not resolved with 200 games (s.e.\ $6$--$13$).



%======================================================================
\section{Prototype}
%======================================================================
% !TEX root = report.tex

\textbf{Current state:} \texttt{wargame\_gui.py} runs the model of \S3 as a \textbf{GUI} (pick each side's plan,
Blue's design and the weather, then step through turns or run many games), as \textbf{batch} runs, or as a \textbf{trace}. \texttt{experiments.py all} regenerates
every number in \S4.

\textbf{How the code reflects the model:} \texttt{Game.step} is the timeline of \S3 with
\eqref{eq:control}--\eqref{eq:payoff}, \texttt{MIPBRPolicy} solves the MIP of \S\ref{sec:mip} each period, and
\texttt{MetaGame} builds $\Gamma$ and runs IESDS, NE enumeration and both LPs. \textbf{Still to build:} learning agents, an attacker that reacts to Blue, Bayesian beliefs and multi-period
policies (\S6).


%======================================================================
\section{Plan for Workshop 2}
%======================================================================
% !TEX root = report.tex
\textbf{Model:} Bayesian beliefs, multi-period best response, sensitivity analysis on $V,h,W,\mu$,
and calendar-based weather from real data. Reformulate the movement budget $F_i$ with considerations for oilers. \textbf{Learning:} a Red agent that learns from repeated play, and other learning
methods on $\Gamma$. Do they converge, drift toward restraint, and exploit rule-based Blue? \textbf{Attacks and defenses:} three Red attacks are coded. A coast guard blockade at SCB and STS blocks
Blue transits, jamming halves Blue's detection in one region, and false contacts appear in Blue's
picture. We will add a Red attacker that reacts to Blue's posture and test randomized
patrols, D1 and D2 as defenses. \textbf{Designs:} D0--D3 and combinations under attack and a long typhoon, and a cost for D3.


\begin{thebibliography}{99}\small\setlength{\itemsep}{0pt}
\bibitem[Green et~al.(2017)]{green2017} Green, M., Hicks, K., Cooper, Z., Schaus, J., and Douglas, J. (2017).
\emph{Countering Coercion in Maritime Asia: The Theory and Practice of Gray Zone Deterrence}. Center for Strategic and International Studies.
\bibitem[Helmus et~al.(2024)]{helmus2024} Helmus, T. C., Romita Grocholski, K., Liggett, T., Rhoades, A. L., Savitz, S., and Palmer, K. (2024).
\emph{Understanding and Countering China's Maritime Gray Zone Operations} (RR-A2954-1). RAND Corporation. \url{https://doi.org/10.7249/RRA2954-1}
\bibitem[Schelling(1966)]{schelling1966} Schelling, T. C. (1966). \emph{Arms and Influence}. Yale University Press.
\bibitem[Shieh et~al.(2012)]{shieh2012} Shieh, E., An, B., Yang, R., Tambe, M., Baldwin, C., DiRenzo, J., Maule, B., and Meyer, G. (2012).
PROTECT: A deployed game theoretic system to protect the ports of the United States. In \emph{Proc.\ AAMAS 2012}, 13--20.
\bibitem[Washburn and Wood(1995)]{washburn1995} Washburn, A., and Wood, K. (1995). Two-person zero-sum games for network interdiction. \emph{Operations Research}, 43(2), 243--251.
\bibitem[Wasser et~al.(2019)]{wasser2019} Wasser, B., Oberholtzer, J., Pettyjohn, S. L., and Mackenzie, W. (2019).
\emph{Gaming Gray Zone Tactics: Design Considerations for a Structured Strategic Game} (RR-2915-A). RAND Corporation. \url{https://doi.org/10.7249/RR2915}
\add{\bibitem[O'Rourke(2025)]{orourke2025} O'Rourke, R. (2025). \emph{China Naval Modernization: Implications for U.S. Navy Capabilities} (RL33153). Congressional Research Service. \url{https://www.congress.gov/crs-product/RL33153}}
\add{\bibitem[Wikipedia(2026a)]{seventhfleet} Wikipedia (2026a). United States Seventh Fleet. \url{https://en.wikipedia.org/wiki/United_States_Seventh_Fleet}, accessed Oct.\ 30, 2026.}
\add{\bibitem[Wikipedia(2026b)]{jmsdf} Wikipedia (2026b). List of active Japan Maritime Self-Defense Force ships. \url{https://en.wikipedia.org/wiki/List_of_active_Japan_Maritime_Self-Defense_Force_ships}, accessed Oct.\ 30, 2026.}
\add{\bibitem[AMTI(2026)]{amti_jcg} Asia Maritime Transparency Initiative. The Japan Coast Guard: resourcing and responsibility. CSIS. \url{https://amti.csis.org/the-japan-coast-guard-resourcing-and-responsibility/}}
\add{\bibitem[Wikipedia(2026c)]{phnavy} Wikipedia (2026c). List of equipment of the Philippine Navy. \url{https://en.wikipedia.org/wiki/List_of_equipment_of_the_Philippine_Navy}, accessed Oct.\ 30, 2026.}
\add{\bibitem[Janes(2026)]{janes_pcg} Janes. Philippine Coast Guard to receive five more OPVs with Japanese assistance. \url{https://janes.com/defence-news/philippine-coast-guard-to-receive-five-more-opvs-with-japanese-assistance}}
\add{\bibitem[Wikipedia(2026d)]{kpn} Wikipedia (2026d). List of active ships of the Korean People's Navy. \url{https://en.wikipedia.org/wiki/List_of_active_ships_of_the_Korean_People\%27s_Navy}, accessed Oct.\ 30, 2026.}
\bibitem[WPNS(2014)]{cues2014} Western Pacific Naval Symposium (2014). \emph{Code for Unplanned Encounters at Sea (CUES)}.
\end{thebibliography}
\noindent\small\textbf{AI usage:} AI models such as Claude and ChatGPT were used to assist in formulating the model, reviewing and
editing the text, writing code, and error checking. The authors reviewed all content and are responsible for it.

\end{document}
